\chapter[Climate Change and the Humanities]{Why Climate Change Is Also a Humanities Question}
\section{A Lump of Coal}
First, pick up a lump of coal.

If your home still has a coal stove, you can actually take one. Otherwise, imagine it: fist-sized, black, lighter than stone, its surface lined like compressed book pages. Hold it and black powder stains your fingers, resisting every attempt to pat it away. Close up, its smell is somewhere between earth and oil.

Here is its history. Around three hundred million years ago, before dinosaurs, hot swamp forests covered much of Earth. Fallen trees soaked in water and were buried under fresh sediment before decaying. Layers accumulated over hundreds of millions of years. Pressure gradually concentrated and stored the forests' carbon, producing the shiny black object in your hand.

What is coal? A flattened forest. More precisely, solid ancient sunlight. Plants stored sunlight three hundred million years ago, were compressed into coal, and were excavated and brought to you. Light it, and the flame releases the sun of that ancient afternoon.

It also releases carbon compressed for those three hundred million years. Entering the air as carbon dioxide, it lets sunlight in while hindering heat's escape. That greenhouse effect can be explained through secondary-school physics.

So far, this is science: chemistry, geology, physics, each statement testable. But our questions begin precisely where that explanation ends. Science calculates emissions from a tonne of coal, not \textbf{who is entitled to burn it}; predicts centimeters of sea-level rise, not \textbf{whose homeland deserves flooding}; describes the next century, not \textbf{who should speak today for its future inhabitants}.

Microscopes cannot answer these; people must. That means reasons, positions, historical accounts, and divided bills: the humanities' work. Hence our question: why is climate change also a humanities problem?

\section{London, December 1952}
Enter a scene with me.

Time: Friday, December 5, 1952. Place: London. Winter arrived early and bitterly cold. Postwar Britain's good coal was exported to repay debts; ordinary people burned cheap, damp, dirty coal producing heavy smoke. That morning, millions of fireplaces and stoves lit together, pouring smoke from chimneys.

Normally smoke rises and disperses. But high pressure stalled overhead, with warmer air above cooler surface air: a temperature inversion. Imagine an invisible pot lid covering the city. Smoke could not rise and thickened near the ground.

For the next five days, London became this:

Yellow fog so thick that streetlights burned by day. Buses stopped, conductors first walking ahead with torches until even that failed. Cinemas closed because audiences could not see screens. An opera in an indoor arena stopped midway as fog entered through doors and performers lost sight of the conductor. At an outlying livestock market, cattle collapsed in groups from suffocation; some survived only with wet sacks over mouths and noses.

The quietest events happened at home. Fog entered doors and windows, and elderly people and those with lung disease struggled to breathe. Deaths over those five days exceeded normal levels by around four thousand. Later studies including following weeks placed the toll potentially above ten thousand. Most victims were elderly, ill, or poor. Poorer residents occupied the foggiest neighborhoods and burned the smokiest coal.

Five years later, Britain passed the Clean Air Act, creating smokeless zones and gradually replacing smoky domestic fuels. London's age as the city of fog finally ended.

Remember three things. First, natural and human disaster cannot be separated here. Nature supplied the inversion and pot lid; humans supplied the coal smoke beneath it. Nature is never humanity's motionless backdrop, but an actor ready to enter. Second, the weakest fell first. Third, people then discussed smoke, coal, deaths, and law, not climate change. Environmental problems become visible through other things: fog, a bill, a funeral, not by stepping into their own spotlight.

\section{Once Facts End, Arguments Begin}
How clear is climate science today? Carbon dioxide's inhibition of heat loss was measured in 1859 by Irish physicist John Tyndall. Swedish chemist Svante Arrhenius calculated warming from coal-driven carbon dioxide increases in 1896. Since 1958, measurements at Hawaii's Mauna Loa have traced a steadily rising atmospheric concentration curve that continues upward. Global average temperature is more than a degree above preindustrial levels; glaciers retreat, seas rise, and extreme weather increases. National academies worldwide have no serious disagreement about these facts.

Where, then, is the problem? Even perfect agreement over the curve removes none of the following arguments, because they concern different questions.

\textbf{First, responsibility.} Everyone helped draw the curve, but pens pressed unequally. Who began first, emitted most, and continues now? Three accounting methods produce three different chief culprits, as we will examine.

\textbf{Second, distribution.} Emission cuts, transition, and sea walls cost money. Who pays? Where do miners work when coal plants close? Who gets new energy jobs? Allocating costs and opportunities is arithmetic plus politics, not physics.

\textbf{Third, time.} Reducing a tonne today chiefly benefits people born seventy or eighty years later. They cannot vote, protest, or criticize online because they do not yet exist. Why spend for them, and why not?

\textbf{Fourth, value.} Is a mangrove forest worth its carbon storage price, its fishing livelihoods, or its own existence? What compensation suffices for a flooded island homeland? Should some things remain unpriced?

Science hands everyone the same examination but cannot answer its hardest questions. Climate difficulty lies not simply in ignorance, but in knowing and \textbf{still facing questions only people can answer}.

Before continuing, test your intuition. An Indian village family buys its first refrigerator; a European family replaces its third car. Both add emissions. Do you already judge them differently? Remember that difference. The chapter asks whether it has grounds.

\section{Islanders, Villagers, and Miners}
Let different people speak, starting with someone losing a homeland.

Two central Pacific countries may be unfamiliar: Tuvalu and Kiribati. They are atoll countries, narrow coral islands around lagoons, averaging roughly two meters above sea level with highest points only a few meters higher. Large storm surges can wash into their centers. Sea-level rise is a headline for us but a surveying problem for them: high tides seep onto capital streets, gardens become too salty for taro, and ancestral drinking wells turn brackish. In 2021, Tuvalu's foreign minister recorded a UN climate-conference address knee-deep in seawater, complete with suit, lectern, and flag. Kiribati bought land in Fiji in 2014, officially for food security, but with an understood possible refuge behind it. Their position: \textbf{our emissions are negligible, yet we move home because of others' smokestacks}.

Hear a different voice from a northern Indian village. Annual per-person electricity use scarcely matches a wealthy household's monthly air-conditioning bill. A headteacher wants fans, a clinic a vaccine refrigerator, young people a workshop. While global newspapers debate reductions, using less energy sounds strange to people who have scarcely begun. India's negotiating position is blunt: you burned coal for two centuries to build factories, railways, and towers, then tell us coal is bad and recommend solar. What happens after sunset? Our coal lies underfoot and people queue for electricity. \textbf{Development is not corruption: later arrivals must live in an environment burdened by those who became rich first.}

Third, hear Germany's Ruhr, with nearly two centuries of coal in its civic identity. Germany's last hard-coal mine closed in 2018 with a formal farewell attended by the president, who received the final lump from the final miner. That dignity was purchased through decades of expensive early retirement, retraining, universities, and new industries. Elsewhere, farewells bring only redundancy notices. Remember a miner's words, in paraphrase: \textbf{for thirty years you called me a builder; overnight my profession became a crime. The furnace is here, but the warmth is in your rooms}.

Now strengthen another side. If you naturally support island states, make the development argument fully enough for its advocates to recognize it.

Which rich-country roads, hospitals, schools, or pensions have no coal or oil behind them? No international rule prohibited their emissions then, and they sent colonies no climate compensation. Once their infrastructure was complete, new rules raised admission prices for everyone else. Those refusing you a ride arrived by vehicle themselves. Meanwhile, rich countries talk transition while retaining much higher per-person emissions: America's remain clearly above China's and far above India's, with several Indians needed to match one American. Requiring hungry people to diet while the well-fed keep ordering is unequal treatment, not environmentalism.

This strong case rests on a simple principle: \textbf{equal rights cannot be monopolized by early arrivals while late arrivals pay}.

But islanders' case also stands, and neither defeats the other. Neither island residents nor Indian villagers burned the coal, yet seas rise regardless. Correct historical accounts cannot halt coral bleaching by moral force. Climate's cruelty is that \textbf{several groups who have done nothing wrong share one shrinking room}.

\section[Sort the Argument onto Four Tables]{Thinking Bridge: Sort the Argument onto Four Tables}
Climate debates at summits or dinner often involve people arguing different questions for half an hour: data, history, money, conscience. Four speakers say climate while fighting four battles.

Use a portable \textbf{four-table sorting method}. Before accepting or rejecting a claim, identify its table.

\textbf{Table one: facts.} What happened, and what will happen? How much warming or sea-level rise? Science arbitrates through data, experiments, and peer review. Evidence matters; allegiance does not. A common indefensible move is losing elsewhere and attacking this table, pretending the thermometer is broken because the bill is unwelcome.

\textbf{Table two: responsibility.} Who caused it? Several ledgers complicate the answer. Who emits most?

\begin{itemize}
\item By \textbf{current annual} emissions, China leads the world.
\item By \textbf{historical accumulation} since the Industrial Revolution, America leads, with Europe and America together accounting for most. Much of today's additional atmospheric carbon came from their two centuries of emissions, persisting for centuries without automatically expiring.
\item By \textbf{per-person} emissions, the answer changes again. America's figure clearly exceeds China's, while India's lies far below the global mean. Dividing China's total among 1.4 billion people immediately reduces it.
\end{itemize}
All three ledgers and answers are real. Arguing responsibility without specifying a ledger resembles fighting with rulers in different units. This table should identify not a single criminal, but \textbf{which ledger you use and why}.

\textbf{Table three: distribution.} What next, who pays, and who gains? Mine closures burden miners; solar and electric vehicles offer new industrial jobs. Carbon taxes raise commuters' fuel costs while opening opportunities for public transport. No arrangement is emissionless; there are \textbf{different distributions of costs and opportunities}. Ask where a policy sends its bill. Charging those least able to escape is often cheapest and least just.

\textbf{Table four: values.} What matters, and whom do we owe? What are future people or wetlands worth? Is preserving a stable smoky job at a distant island's expense wrong? Here there is no referee, only argument. Offer reasons and submit them to others' reasons.

The rule is \textbf{sort before arguing}. Saying developing countries must immediately reduce coal is true at the facts table, incomplete at responsibility because history remains uncounted, negotiable at distribution because someone must finance grids, and a position at values because present poverty and future seas must be weighed. One sentence occupies four tables; rebutting it with just one table's rules misses. Television and trending comments miss this way daily.

\section{Common but Differentiated in a Convention}
Read a primary source from a document almost every country has signed.

States meeting in Rio de Janeiro in 1992 adopted the UN Framework Convention on Climate Change. Article 3, paragraph 1 says, in broad translation:

\begin{quote}
Parties should protect the climate system for present and future generations on the basis of equity, according to their common but differentiated responsibilities and respective capabilities. Developed-country parties should therefore lead in addressing climate change and its adverse effects.
\end{quote}
Focus on \textbf{common but differentiated responsibilities}. These seven characters embodied the broadest agreement then possible and condensed thirty subsequent years of argument.

Common means one atmosphere receives everyone's emissions, leaving nobody outside. Differentiated invokes the ledgers: histories, capacities, and wealth differ, so bills cannot be split equally. Developed countries leading is plainer still: earlier emitters cut first, wealthier countries pay first.

Why this wording? It is a \textbf{carefully designed compromise}. Islanders wanted immediate, large, binding cuts; developing countries wanted historical recognition without blocked development; wealthy states wanted shared action rather than exclusive blame. None persuaded the others, so each took half a sentence, assembled into something all could sign. One diplomatic secret is that \textbf{signable sentences often let every party read what it wants inside}.

Subsequent history alternates emphasis. The 1997 Kyoto Protocol implemented differentiation: only developed countries received legally binding reduction targets, not developing ones. America ultimately did not ratify, invoking common responsibility against exclusive constraint. The 2015 Paris Agreement moved toward common action: countries submit and update their own plans under worldwide scrutiny instead of receiving imposed targets. This may be more realistic, placing America, China, and India in one framework, or weaker, with self-assigned homework and self-checking. Both readings find something real in the convention.

This is why primary sources matter. Those seven characters are not a sage's answer but an arena. They recognize \textbf{historical emissions} through differentiation, \textbf{development rights} through respective capabilities, and \textbf{future generations} as stakeholders. They settle none of them. Thirty years of negotiations essentially ask how differentiated responsibilities should be.

\section[Easter Island: Three Stories]{One Island, Three Stories: Easter Island Revisited}
Now compare genuinely different explanations of one frequently repeated climate parable.

First, facts. More than three thousand kilometers from the nearest continent lies Rapa Nui, outsiders' Easter Island. Its Polynesian ancestors erected hundreds of immense figures, the tallest nearly ten meters. When Europeans arrived in 1722, almost no large trees remained. How did an island of volcanic rock, grass, and statues produce those statues? \textbf{What happened here?}

\textbf{Explanation one: ecological suicide.} Jared Diamond's Collapse systematically presents the popular account: population growth and demand for statue transport, boats, and fuel eliminated palms. Deforestation eroded soils, prevented fishing-boat construction, undermined food supply, and produced warfare and population collapse. The tale became a standard warning of humanity consuming its homeland. Its support: forest loss is real. Its limit: who caused that loss and how severely islanders suffered are less firmly established than the narrative suggests.

\textbf{Explanation two: other culprits, more capable islanders.} Researchers including Terry Hunt and Carl Lipo offer another picture. Polynesian rats arriving with settlers ate palm seeds and may have substantially prevented regeneration. Islanders adapted through rock gardens, conserving moisture and nutrients to grow sufficient sweet potatoes in poor soil. Experiments also revisited statue movement: around a dozen people with ropes could rock statues into walking, without tree-trunk rollers. Supporting evidence includes gnawed ancient seeds and field experiments. Its limit: successful adaptation does not automatically prove no decline.

\textbf{Explanation three: collapse followed European arrival.} Historical records describe Peruvian slavers repeatedly raiding from 1862, taking over a thousand islanders for labor in guano mines. Few returned alive, bringing smallpox. By 1877, only a little over a hundred residents remained. The dramatic population cliff likely followed outside contact, with slavers and disease among its culprits. Blaming foolish tree-cutters transfers responsibility from outsiders to victims. Ships' and church records support raids and epidemics. This account does not deny deforestation; it rejects blaming islanders for every disaster.

Each explanation holds evidence without exhausting truth. Narrative selection matters: \textbf{which Easter Island you tell depends on the conclusion you want today's audience to draw}.

Face this directly: \textbf{disaster and hope narratives are powerful political tools, each with costs}. Ecological suicide warns us not to repeat a mistake. But quietly deleting slavers inflicts a second injury on the dead: first population stolen, then reputation. Hope can mislead too. Praising ingenuity and technological solutions alone invites complacency about emissions. Disaster can wake people or paralyze them into inevitable-doom resignation. Hope can sustain action or anesthetize people into waiting for rescue. Storytellers hold steering wheels, so narratives require inspection, including this book's.

This island also provides an \textbf{easily misread counterexample}. Environmental decline causing civilizational collapse seems obvious until its most famous case produces divergent evidence. That does not make environmental danger unreal. It means \textbf{the more pleasing the story, the more insistently you should ask for evidence}.

\section[What Are Future People Worth Today?]{Deeper Challenge: What Are Future People Worth Today?}
Our next tool looks economic but contains an entire ethics.

Imagine I owe you a hundred yuan, then announce repayment in eighty years. You feel cheated: that future hundred is not today's hundred. Economists formalize this intuition as \textbf{discounting}, reducing future money or benefits to present value. Greater discounts make the future worth less. The conversion rate is the \textbf{discount rate}.

Ordinary accounting uses it readily. Climate exposes its teeth, because present costs purchase benefits chiefly seventy or eighty years later. What is preventing a flood eighty years hence worth now?

Calculate a simple loss of one hundred yuan eighty years from today.

\begin{itemize}
\item At five percent annual discounting, its present value is roughly \textbf{two yuan}. Spending more than two today on prevention therefore appears uneconomic.
\item At 1.4 percent, the same loss is worth around \textbf{thirty-three yuan} today, immediately deserving more serious attention.
\end{itemize}
Two and thirty-three differ sixteenfold, enough to influence coal closures, heavy fuel taxes, or hundreds of billions for clean energy. \textbf{The deciding factor is a percentage, not an experiment or dataset.} Where does that percentage come from?

The British government's 2006 Stern Review called climate change history's greatest market failure and argued that spending roughly one percent of annual global output on strong reductions would cost far less than later damage. Its low discount rate resembles the thirty-three-yuan case. Another position, associated with American economist and later Nobel laureate William Nordhaus, anchors discounting in actual investment returns. Higher rates, nearer the two-yuan case, reverse the emphasis: reduce emissions gradually, leaving richer descendants more of the cost.

Apparently economists dispute a technical parameter. Underneath lies \textbf{ethics dressed as mathematics}.

A high rate argues future people will probably be richer through progress, so transfers from rich to richer are less urgent, while present investment elsewhere can earn returns. A low rate argues \textbf{someone's life should not automatically lose value merely because they are born later}. We would not halve an infant's worth for living next door across a border; why do so for living a century later?

This is \textbf{intergenerational justice}. Future people cannot vote, hire lawyers, or refuse, yet are the largest stakeholder group, inheriting atmospheric carbon for centuries. Denying rights because they do not yet exist would imply no wrong in planting a bomb to explode eighty years later because its victims are unborn. Nobody accepts that conclusion. Equal rights, however, raise conflicts between today's poor needing heat and tomorrow's coastlines. The convention names present and future interests together without settling their collisions.

Deeper still, \textbf{the calculation assumes everything can become money}. What compensates a flooded homeland, an unspoken language, vanished mangroves, or a bird surviving only in legend? Some call pricing nature its only route into decisions, where unpriced means zero. Others call it violence: not everything inhabits markets. Discount-rate opponents share convertibility as a premise. The fourth, values table remains for things \textbf{refusing conversion}.

Recognize the tool's boundary. It can assess economic worth, not establish justice. Mathematics handles arithmetic; conscience handles accounts that cannot be calculated.

\section{Climate Negotiations Downstairs}
Return to today, where the issue sits close enough to escape recognition.

First, nature replies. In July 2021, Zhengzhou experienced extraordinary rainfall. An official investigation reported 398 deaths and missing people across Henan, including commuters trapped in a subway carriage. In summer 2022, the Yangtze basin suffered one of its strongest recorded heatwaves. Hydropower-dependent Sichuan lost generating capacity, factories faced power limits, and shops dimmed lights. Like London's fog, these show \textbf{nature is never scenery}. Cities built around an assumption of quiet nature discover they have no lines when it speaks. Treating it as fixed backdrop is an expensive human-centered illusion, with bills arriving.

Second, transition distributes costs and chances. News shows mining cities where closures obsolete a generation's skills alongside new solar and electric-vehicle industries. Within little more than a decade, China developed world-leading battery and photovoltaic sectors and jobs in different cities. One transition is a cliff for some and an elevator for others, often in different provinces. Thus clean-energy arguments persist: \textbf{one speaker stands in an elevator, another beside a cliff}. The Ruhr's long, costly farewell teaches not merely admiration for civilization, but that \textbf{a just transition requires building bridges for those at the cliff, and bridges cost money}. Financing them remains an agonizing budget issue everywhere.

Third, examine your own table: disposable cutlery with deliveries, phone replacement frequency, classroom air-conditioning temperature, and fast-fashion use. These choices matter and add up. But keep proportions honest. Power, industry, and transport dominate emissions. Reducing the entire climate issue to children's plastic straws shifts responsibility: \textbf{structural bills become individual moral homework}. Personal choices are not meaningless, but their strongest effect may be voting, speaking, and redefining normal life rather than saving a few grams. When enough people reject replacing phones three times yearly as normal, factory schedules can change.

Narrative conflict also lives downstairs. Platforms offer catastrophe---collapsing glaciers, ten-year deadlines, too late---until you feel numb, and bland hope---innovation, green living, a promising future---until you feel equally numb. Both lead to inaction. What does defensible hope look like? Consider Saihanba in Hebei, near Beijing. Since 1962, three generations of foresters have planted over a million mu of forest on once-barren sandlands, producing a forest capable of altering local climate. Unlike technology-will-fix-everything slogans, \textbf{this is a receipt, not reassurance}. Disaster identifies the fire; hope needs receipts to avoid becoming an empty check. Distinguishing them is basic online literacy.

Finally, hear Mencius advising King Hui of Liang more than two thousand years ago, in ``Liang Hui Wang I'':

\begin{quote}
Do not disrupt farming seasons, and grain will be more than can be eaten. Keep fine-meshed nets from deep pools, and fish and turtles will be more than can be eaten. Send axes into forests in season, and timber will be more than can be used.
\end{quote}
The point is timely farming, avoiding fine-meshed fishing, and seasonal cutting. Mencius is not yet a modern environmentalist; abundant usable timber remains his concern, with nature as human resource. Yet he recognizes something enduring: \textbf{nature has rhythms, and human taking must leave room for them}. In season is agricultural civilization's early agreement with nature: take, but do not dictate the timetable. Today that contract extends from one forest to the whole atmosphere.

\section{For You: What Do You Owe Strangers?}
Return to the coal. Burned, it releases ancient sunlight while carbon remains in everyone's sky. Answer the following with reasons.

Imagine two strangers.

The first is a child your age on a Kiribati atoll. You will never meet, share no language, and may forget the country after a geography exam. Rising seas salt the family's garden. The child has burned little coal, while your charging phone, cooled classroom, and delivered parcels add water to that encroaching sea.

The second is farther away: a child a century hence, not yet existent, unable to be photographed making a seawater speech. Yet your ledgers connect. Today's carbon waits overhead at their birth; today's trees, sea walls, and solar plants become their inheritance too.

\textbf{How much do you owe these strangers?}

More concretely, would you accept thirty extra yuan on the family's monthly electricity bill to preserve that island another fifty years? Would your generation, responsible for few emissions so far, surrender easy conveniences for an unborn child? Do you reason that you owe them, that historical emitters should pay first, or that your own life has scarcely begun and should not start with settlement of a bill?

Weigh the other direction too. The Indian child awaiting a fan is a stranger, as is the thirty-year miner. Every climate-ledger page carries strangers' names; defending one page may send its bill to another.

There is no standard answer, but one minimum requirement: identify your table. Are you presenting facts, allocating responsibility, distributing costs, or stating conscience? And which ledger records your own place for future people?

One day, you will be that future person.
